\documentclass[10pt,a4paper]{amsart}
\usepackage[margin=19mm]{geometry}
\usepackage{amsmath,amssymb,mathtools}
\usepackage[scr=rsfs,cal=euler]{mathalfa}
\usepackage{xfrac}
\usepackage[unicode,hidelinks,bookmarks=false]{hyperref}
\newcommand{\ie}{i.e.\ }
\newcommand{\cf}{cf.\ }
\newcommand{\resp}{respectively\ }
\newcommand{\Subsection}{\subsection}
\providecommand{\nobreakdash}{\nobreak\hbox{-}}
\setlength{\emergencystretch}{2em}
\multlinegap0pt
\usepackage{bm}
\usepackage{bbm}
\usepackage{dsfont}
\usepackage{xspace}
\usepackage{cancel}
\usepackage{mathtools}
\usepackage{cleveref}
\usepackage{amsthm}
\makeatletter
\@ifundefined{Theorem}{\newtheorem{Theorem}{Theorem}}{}
\@ifundefined{Lemma}{\newtheorem{Lemma}[Theorem]{Lemma}}{}
\@ifundefined{Remark}{\newtheorem{Remark}[Theorem]{Remark}}{}
\makeatother
\newcommand{\D}{\delta}
\newcommand{\R}{\mathbb{R}}
\title[Polyhedral results for two classes of submodular sets with G]{Polyhedral results for two classes of submodular sets with GUB constraints}
\author{Weikang Qian, Keyan Li, Wei-Kun Chen, Yu-Hong Dai}
\date{Source: arXiv:2601.18360v1 (CC BY 4.0)}
\begin{document}
\maketitle

\noindent\emph{Source and licence.} Polyhedral results for two classes of submodular sets with GUB constraints. Version arXiv:2601.18360v1 (CC BY 4.0), distributed under the Creative Commons Attribution 4.0 International licence, as recorded in the arXiv licence metadata for this exact version and shown on the abstract page https://arxiv.org/abs/2601.18360v1. Full rights notice: licenses/arxiv-2601.18360-CC-BY-4.0.txt.\par\smallskip

\noindent\textit{Editorial scope.} Counted unit: the Lemma whose source label is \texttt{lem:induction} in the article (source0.tex lines 620--622, source label \texttt{lem:induction}), delivered together with the article's own complete original proof of it (source0.tex lines 623--725, one \texttt{proof} environment of 103 lines). No mathematics is written, repaired or re-derived: statement and proof are the article's own text line for line. Required context retained uncounted: asumption:a (lines 308--310); prob:equlift (lines 498--502); lem:propertyU (lines 513--516); problem:defh (lines 518--527); lem:propertyh (lines 535--537); def:W (lines 543--545); lem:feasibilityW (lines 554--556); induction hypothesis 1 (lines 573--575); lem:validity (lines 577--579); lem:propertyeta (lines 583--588); property:eta (lines 585--587). Every definition, display and label of those lines is retained unchanged. Omitted from this package and not redistributed: 1--307, 311--497, 503--512, 517--517, 528--534, 538--542, 546--553, 557--572, 576--576, 580--582, 588--619, 726--1745. Nothing inside a retained excerpt is omitted or altered: every retained body is delivered exactly as the article's lines are written, so no omission receipt of any kind accompanies this package. Citations: the retained excerpts carry no citation command at all, so no bibliography entry of the article is transcribed and no reference is redistributed. Reference adaptation: none was needed and none was made; every reference inside the delivered unit resolves inside it, so no unresolved reference or citation is produced. Printed slips kept literal: no printed word, symbol, hypothesis or number of the counted statement or of its proof was altered. Layout and class: the article's own class is \texttt{article} with packages including geometry, algcompatible, caption, titlesec, hyperref, authblk, graphicx, multirow, amsmath, amssymb, amsfonts, amsthm, mathrsfs, appendix; the wrapper is the lane's pinned \texttt{amsart} editorial wrapper at 10pt on A4, which restates the theorem-like environments the retained text needs and adds the article's own macro definitions that the retained text needs (\texttt{\textbackslash D}, \texttt{\textbackslash R}), with extra wrapper packages (bm, bbm, dsfont, xspace, cancel, mathtools, cleveref). The delivered numbering is the unit's own. Assets: no retained span contains a figure environment, a graphics-inclusion command, a PDF-page inclusion, a table or a TikZ picture; no image file is copied into this package, the package declares no asset dependency and no image file is redistributed with it. Licence: version arXiv:2601.18360v1 of this article is distributed under the Creative Commons Attribution 4.0 International licence, as recorded in the arXiv licence metadata for this exact version and shown on the abstract page https://arxiv.org/abs/2601.18360v1. A full rights notice is delivered at licenses/arxiv-2601.18360-CC-BY-4.0.txt. Characters: every retained excerpt is pure ASCII, so the delivered engine, XeLaTeX with its default Latin Modern text font, reports no missing character.
\begin{equation}\label{asumption:a}
	a_{i_{k-1}+1}\leq a_{i_{k-1}+2}\leq\cdots\leq a_{i_{k-1}+|N_k|},~\forall~ k\in[t].
\end{equation}

\begin{equation}\label{prob:equlift}
	\eta^\D_{\D_j} = \min_S \left\{  f(a(S))-\sum_{i\in S \backslash\D_j}\eta^{\D}_i \,:\,
	S \subseteq \D(j), ~ |S \cap N_k |\leq 1, ~\forall~k\in[t], ~\D_j \in S
	 \right\}.
\end{equation}

\begin{Remark}\label{lem:propertyU}
	For any $j \in [n]$, (i) $|U_j^\delta \cup N_k|\leq 1$ holds for all $k \in [t]$; and 
	(ii) $a(U^{\D}_j)\geq a(S)$ holds for all $S \subseteq \D(j)$ with $|S \cap N_k |\leq 1$ for $k\in[t]$.
\end{Remark}

\begin{equation}\label{problem:defh}
	h^{\D}_j:=
	\begin{cases}
		j, & 
		~\text{if}~ \D_j= \max_{i\in N_k\cap \D(j)} i,\\
	    \min\{\ell\in [j-1]\,:\,\D_{\ell}>\D_j,
	    ~\D_{\ell}\in N_k\}, 
	    & ~\text{if}~ \D_j < \max_{i\in N_k \cap \D(j)} i.
	\end{cases}
\end{equation}

\begin{Lemma}\label{lem:propertyh}
	$\D_{h^\D_j}\in U^\D_{h^\D_j}$ holds for all $j\in[n]$.
\end{Lemma}

	\begin{equation}\label{def:W}
		W^\D_j:=U^\D_{ h^\D_j}\backslash \D_{h^\D_j}\cup \D_j
	\end{equation}

\begin{Lemma}\label{lem:feasibilityW}
	For $j\in[n]$, $W_j^\D$ is a feasible solution of  problem \eqref{prob:equlift}.
\end{Lemma}

\begin{equation}\label{induction hypothesis 1}
	\eta^{\D}_{\D_j}=f(a(W_j^\D))-\sum_{i\in W_j^\D \backslash\D_j}\eta^{\D}_i, ~\forall~ j\in[r-1].
\end{equation}

\begin{Lemma}\label{lem:validity}
	For any set $S \subseteq [n]$ satisfying $|S \cap N_k |\leq 1$ for $k\in[t]$, the inequality $f(a(S))\geq \sum_{i\in S}\eta^\D_i$ holds.
\end{Lemma}

\begin{Lemma}\label{lem:propertyeta}
	Suppose that \eqref{induction hypothesis 1} holds. Then it follows that
	\begin{equation}\label{property:eta}
		f(a(U^\D_j))=\sum_{i\in U^{\D}_j}\eta^{\D}_i, ~\forall~ j\in[r-1].
	\end{equation}
\end{Lemma}

\begin{Lemma}\label{lem:induction}
	Suppose that the induction hypothesis \eqref{induction hypothesis 1} holds. Then $W^{\D}_r$ is an optimal solution of problem \eqref{prob:equlift} with $j=r$.
\end{Lemma}

\begin{proof}
	Letting $S$ be an arbitrary feasible solution of the lifting problem \eqref{prob:equlift} with $j=r$, then 
	$\D_r\in S$, $S\subseteq\D(r)$, and $|S \cap N_k |\leq 1$ for $k\in[t]$.
	By \cref{lem:feasibilityW}, it suffices to show $f(a(W^\D_r))-\sum_{i\in W^{\D}_r\backslash \D_r} \eta^\D_i \leq f(a(S))-\sum_{i\in S\backslash \D_r}\eta^\D_i$, 
	or equivalently, 
	\begin{equation}\label{ourgoal2}
		 f(a(S))-f(a(W^\D_r)) \geq \sum_{i\in S\backslash \D_r}\eta^\D_i-\sum_{i\in W^{\D}_r\backslash \D_r} \eta^\D_i.
	\end{equation}
	Let $k' \in [t]$ be such that  $\D_r\in N_{k'}$.
    We next consider the following two cases (i) $a(W^{\D}_r)<a(S)$ and (ii) $a(W^{\D}_r)\geq a(S)$, separately. \\[8pt]
	{(i)} $a(W^{\D}_r)<a(S)$.
 	If $h^\D_r=r$, then $W^\D_r=U^\D_r$, which, 
	together with \cref{lem:propertyU}(ii), implies $a(S)\leq a(U^{\D}_r)=a(W^{\D}_r)$,
	a contradiction.
	Thus, by the definition of $h^\D_r$ in \eqref{problem:defh}, $h^\D_r< r$ and $\D_{h^\D_r}> \D_{r}$ must hold.
	Observe that
	\begin{equation}\label{tmpeq1}
		\begin{aligned}
			 &\sum_{i\in S\backslash\D_r}
			\eta^{\D}_i-\sum_{i\in W^{\D}_{r}\backslash\D_{r}}\eta^{\D}_i
			 \stackrel{(a)}{=} \sum_{i\in S\backslash \D_r}
			\eta^{\D}_i-\sum_{i\in U^{\D}_{h^\D_r}\backslash\D_{h^\D_r}}\eta^{\D}_i \\
			&\qquad \stackrel{(b)}{=} \sum_{i\in S\backslash \D_r\cup\D_{h^\D_r}}
			\eta^{\D}_i-\sum_{i\in U^{\D}_{h^\D_r}}\eta^{\D}_i 
			\stackrel{(c)}{=}  \sum_{i\in S\backslash \D_r\cup\D_{h^\D_r}}\eta^{\D}_i-f(a(U^{\D}_{h^\D_r})),
		\end{aligned}
	\end{equation}
	where (a) follows from the definition of $W^\D_j$ in \eqref{def:W} and $\D_r\notin U^\D_{h^\D_r}\backslash \D_{h^\D_r}$;
	(b) follows from $\D_{h^\D_r}\notin S\backslash \D_r$ (by $\D_r,\D_{h^\D_r}\in N_{k'}$, $|S \cap N_{k'}|\leq 1$, and $\D_r \in S$) and $\D_{h^\D_r}\in U^{\D}_{h^\D_r}$ (by \cref{lem:propertyh});
	and (c) follows from the induction hypothesis in \eqref{induction hypothesis 1} and \cref{lem:propertyeta}; that is, \eqref{property:eta} with $j=h^\D_r<r$ holds.
	 On the other hand,
	 \begin{equation}\label{tmpeq2}
	 	\begin{aligned}
	 		&f(a(S))-f(a(W^{\D}_r)) \stackrel{(a)}{=} f(a(S\backslash \D_r)+a_{\D_r})
	 		-f(a(U^{\D}_{h^\D_r}\backslash \D_{h^\D_r})+a_{\D_r})\\
	 		&\qquad \stackrel{(b)}{\geq} f(a(S\backslash \D_r)+a_{\D_{h^\D_r}})
	 		-f(a(U^{\D}_{h^\D_r}\backslash \D_{h^\D_r})+a_{\D_{h^\D_r}}) \\
	 		&\qquad \stackrel{(c)}{=}f(a(S\backslash \D_r\cup\D_{h^\D_r}))
	 		-f(a(U^{\D}_{h^\D_r}))\stackrel{(d)}{\geq} \sum_{i\in S\backslash \D_r\cup\D_{h^\D_r}}
	 		\eta^{\D}_i-f(a(U^{\D}_{h^\D_r})),
	 	\end{aligned}
	 \end{equation}
	 where (a) follows from $\D_r\in S$ and $\D_r\notin U^\D_{h^\D_r}\backslash \D_{h^\D_r}$; 
	 (b) follows from  $a(S\backslash \D_r)= a(S)-a_{\D_r}> 
	 a(W^{\D}_r)-a_{\D_r}= a(U^{\D}_{h^\D_r}\backslash\D_{h^\D_r})$,
	 $a_{\D_{h^\D_r}}\geq a_{\D_r}$ (by $\D_{h^\D_r},\D_r\in N_{k'}$, $\D_{h^\D_r}> \D_{r}$, 
	 and \eqref{asumption:a}), and the concavity of $f$ (that is, $f(y_2)-f(y_1)\geq f(y_2+d)-f(y_1+d)$ for any $d\in\R_+$ and $y_1,y_2\in\R$ with $y_1\leq y_2$);
	 (c) follows from $\D_{h^\D_r}\notin S\backslash \D_r$ 
	 (by $\D_r,\D_{h^\D_r}\in N_{k'}$ and $|S \cap N_{k'}|\leq 1$)
	 and $\D_{h^\D_r}\in U^{\D}_{h^\D_r}$ (by \cref{lem:propertyh});
	 and (d) follows from \cref{lem:validity} and $|(S\backslash \D_r\cup \D_{h^\D_r})\cap N_k|\leq 1$ for $k \in [t]$.
	 Combining \eqref{tmpeq1} and \eqref{tmpeq2} yields the desired result in \eqref{ourgoal2}.\\[8pt]
	{(ii)} $a(W^{\D}_r)\geq a(S)$.
	Then, by $W_r^\D= U^\D_{ h^\D_r}\backslash \D_{h^\D_r}\cup \D_r$, $\D_r \notin U^\D_{ h^\D_r}\backslash \D_{h^\D_r}$, and $\D_r \in S$, it follows
	\begin{equation}\label{tmpeq3}
		a(U^\D_{h^\D_r}\backslash\D_{h^\D_r})
		= a(W^\D_r)-a_{\D_r}\geq a(S)-a_{\D_r}= a(S\backslash \D_r).
	\end{equation}
	Let
	\begin{equation}\label{defU'}
		U':=
		\begin{cases}
			\varnothing & 
			~\text{if}~  N_{k'}\cap\D(h^\D_r-1)=\varnothing,\\
			\{\max_{i\in N_{k'}\cap\D(h^\D_r-1)}i\},
			& ~\text{otherwise}.
		\end{cases}
	\end{equation}
	By definition, it follows that
	\begin{equation}\label{J'-1J'}
		U^\D_{h^\D_r-1}\backslash U'= \bigcup_{k\in [t]\backslash k'}
		\left\{\max_{i\in N_k\cap\D(h^\D_r-1)}i\right\}= \bigcup_{k\in [t]\backslash k'}
		\left\{\max_{i\in N_k\cap\D(h^\D_r)}i\right\}=U^\D_{h^\D_r}\backslash \D_{h^\D_r}
	\end{equation}
	and hence
	\begin{equation}\label{tmpeq4}
		a(U^{\D}_{h^\D_r-1})= a(U^{\D}_{h^\D_r}\backslash\D_{h^\D_r})+a(U').
	\end{equation}
	From the definition of $h^\D_r$ in \eqref{problem:defh}, 
	$\D_{\ell}<\D_r$ holds for all $\ell<h^\D_r$ with $\D_{\ell}\in N_{k'}$.
	This, together with \eqref{asumption:a} and the definition of $U'$ in \eqref{defU'}, implies $a(U')\leq a_{\D_r}$.
	To prove the desired result in \eqref{ourgoal2}, observe that
	\begin{equation*}
		\begin{aligned}
			&f(a(S))-f(a(W^{\D}_r))=-[f(a(U^{\D}_{h^\D_r}
			\backslash \D_{h^\D_r})+a_{\D_r})-f(a(S\backslash \D_r)+a_{\D_r})]\\
			&\qquad \stackrel{(a)}{\geq} -[f(a(U^\D_{h^\D_r}\backslash \D_{h^\D_r})+a(U'))-f(a(S\backslash \D_r)+a(U'))]
			\stackrel{(b)}{=}f(a(S\backslash \D_r\cup U'))-f(a(U^{\D}_{h^\D_r-1}))\\
			&\qquad \stackrel{(c)}{\geq} \sum_{i\in S\backslash \D_r\cup U'}\eta^{\D}_i-f(a(U^{\D}_{h^\D_r-1}))
			\stackrel{(d)}{=}\sum_{i\in S\backslash \D_r\cup U'}\eta^{\D}_i- \sum_{i\in U^{\D}_{h^\D_r-1}}\eta^{\D}_i\\
			&\qquad \stackrel{(e)}{=} 
			\sum_{i\in S\backslash \D_r}\eta^{\D}_i-\sum_{i\in U^{\D}_{h^\D_r-1}\backslash U'}\eta^{\D}_i
			\stackrel{(f)}{=} \sum_{i\in S\backslash\D_r}\eta^{\D}_i-\sum_{i\in W^{\D}_{r}\backslash\D_{r}}
			\eta^{\D}_i,
		\end{aligned}
	\end{equation*} 
	where (a) follows from \eqref{tmpeq3}, $a(U')\leq a_{\D_r}$, and the concavity of $f$;
	(b) follows from \eqref{tmpeq4};
	(c) follows from \cref{lem:validity} and $|(S\backslash \D_r\cup U')\cap N_k|\leq 1$ for $k \in [t]$;
	(d) follows from the induction hypothesis in \eqref{induction hypothesis 1} and \cref{lem:propertyeta}, that is, \eqref{property:eta} with $j=h^\D_r-1\leq r-1$ holds;
	(e) follows from $U'\cap (S\backslash \D_r)=\varnothing$ and $U'\subseteq U^{\D}_{h^\D_r-1}$; and 
	(f) follows from \eqref{J'-1J'} and $U^\D_{h^\D_r}\backslash \D_{h^\D_r}=W^\D_r\backslash \D_r$. 
\end{proof}

\end{document}
